\section{Appendix A. Mathematical Derivations}

\subsection{A.1 Recurrence Unrolling}

For the HGRN-2 recurrence:

\[
s\_t = f\_t \odot s\_{t-1} + g\_t
\tag{A.1}
\]

with $s_0 = 0$, the closed form is:

\[
s\_t = \sum\_{j=1}^{t} \left( \prod\_{i=j+1}^{t} f\_i \right) g\_j
\tag{A.2}
\]

The effective decay from position $j$ to position $t$ is:

\[
D\_{t, j} = \prod\_{i=j+1}^{t} f\_i = \exp\left(\sum\_{i=j+1}^{t} \log f\_i\right)
\tag{A.3}
\]

\subsection{A.2 Chunked Scan}

For a chunk $[s, e)$ with $C = e - s$ timesteps, define:

\[
A\_k = \sum\_{i=s}^{k} \log f\_i, \quad k \in [s, e)
\tag{A.4}
\]

Then within the chunk:

\[
h\_k^\text{intra} = \sum\_{j=s}^{k} \exp(A\_k - A\_j) \cdot g\_j
\tag{A.5}
\]

Combined with previous chunk's final state $h_\text{prev}$:

\[
h\_k^\text{chunk} = h\_\text{prev} \cdot \exp(A\_k) + h\_k^\text{intra}
\tag{A.6}
\]

This is exact (no approximation).

\subsection{A.3 Bounded Forget Gate}

The forget gate:

\[
f\_t = \exp(-\text{softplus}(W\_f x\_t))
\tag{A.7}
\]

Softplus is always positive: $\text{softplus}(u) > 0$. So $f_t \in (0, 1)$.

\textbf{Clamp in log-space:}

\[
\log f\_t \leftarrow \text{clamp}(- \text{softplus}(u), -3, -10^{-4})
\tag{A.8}
\]

So $f_t \in [e^{-3}, e^{-10^{-4}}] \approx [0.05, 1]$.

\textbf{Why clamp?} Without clamping, $\text{softplus}(u)$ can be very large (if 
$u$ is large), giving $f \approx 0$ (no memory). Or very small, giving 
$f \approx 1$ (no update). The clamp prevents both.

\subsection{A.4 Bidirectional Combination}

For encoders:

\[
h\_t = h\_t^\text{fwd} + h\_t^\text{bwd}
\tag{A.9}
\]

\textbf{Why sum, not concatenate?} Sum preserves dimension. Concatenate would 
double the hidden size. Averaging would halve the magnitude. Sum is the 
standard choice in BiRNN implementations.

\subsection{A.5 LoRA Recovery}

Base weights are frozen. LoRA adapters:

\[
\Delta W\_v = \frac{\alpha}{r} B\_v A\_v, \quad
\Delta W\_o = \frac{\alpha}{r} B\_o A\_o
\tag{A.10}
\]

with $B_v, B_o = 0$ initially.

\subsection{A.6 Complexity}

For a single layer of dim $d$, sequence $L$:

\begin{table}[H]
\centering
\small
\begin{tabular}{lll}
\toprule
\textbf{Component} & \textbf{Time} & \textbf{Memory} \\
\midrule
$W_f x$ & $L d^2$ & $L d$ \\
$W_i x$ & $L d^2$ & $L d$ \\
$W_v x$ & $L d^2$ & $L d$ \\
$W_g x$ & $L d^2$ & $L d$ \\
Chunked scan & $L d$ & $C^2 + L d$ \\
$W_o$ & $L d^2$ & $L d$ \\
\textbf{Total} & $5 L d^2 + L d$ & $\mathcal{O}(L d)$ \\
\bottomrule
\end{tabular}
\end{table}

\subsection{A.7 Chunk Size and Parallelism}

For chunk size $C$, the number of sequential steps is $L / C$. The 
intra-chunk work is $C^2$ per chunk, total $L \cdot C$. So:

\begin{itemize}
\item \textbf{Larger $C$:} more parallel work per chunk, fewer sequential steps, but
  $C^2$ memory and compute per chunk.
\item \textbf{Smaller $C$:} less memory, more sequential steps.

\end{itemize}
For L = 512, C = 32: 16 sequential chunks. This is a good balance.
\begin{lstlisting}

---
\end{lstlisting}